\documentclass{article}
\usepackage{graphicx} % Required for inserting images

\title{CS 345 Draft Proposal}
\author{Abdulrahman Alshahrani (KAUST ID: 182856)}
\date{February 2026}

\begin{document}

\maketitle

\section{Project Objective}

In this reproduction project, I aim to reproduce the Agentic Context Engineering (ACE) framework experiments under MAESTRO's framework, utilizing MAESTRO's system traces to provide further insight into ACE's CPU, memory, and communication overhead. Additionally, I aim to investigate if using better (and more expensive) LLMs for the Reflector agent can provide a better cost-accuracy result, an aspect identified by ACE authors but not explored.

\section{Background}

\subsection{Summary of the Agentic Context Engineering paper}

The Agentic Context Engineering (ACE) paper introduces a new framework of managing agent context effectively, aiming to support the development of long-living and self-improving LLM agents. The framework addresses the major flaw in current methods of managing context in long workflows, named as context collapse and brevity bias, where domain-specific and important information gets omitted over long sessions. They resolve this issue by treating contexts as evolving playbooks, utilizing a three-agent workflow to generate, reflect, and curate better instructions overtime. Their proposed framework produces the best results on the AppWorld and FiNER benchmarks, against similar approaches.

\subsection{Summary of the MAESTRO paper}

The MAESTRO paper provides a standardized framework to allow for the testing, observability, and measure reliability of multi-agent systems. Beyond the simple task-success benchmarks that already exist, MAESTRO offers system-level traces of model interaction with the system, allowing for the analysis of CPU, memory, and network usage, as well as specifying metrics to measure inter-run consistency of call-graphs.

\section{Technical Idea}

- Implement ACE's framework using the adapters provided by MAESTRO to get a standard telemetry interface.
\\
- Using MAESTRO, we will investigate when the ACE framework run the three-model workflow, and analyze the CPU, memory, and network usage.
\\
- In one settings, we would replace ACE's Reflector agent with a better LLM (e.g. GPT-5.2) whereas the other agents would use a smaller, cheaper LLM (e.g. GPT-5-nano). With this, we would analyze the resource usage and the overall accuracy differences.

\section{Evaluation Metrics}

- As specified above, resource footprint will be measured with MAESTRO traces (CPU, Memory, and network usage).
\\
- We can also measure the consistency of call graphs, using the same Jaccard Similarity metric used in the MAESTRO paper.

\section{Benchmarks to be re-run}

- We will focus on the AppWorld and FiNER benchmarks mentioned in the ACE paper.

\section{Resources}

- Agentic Context Engineering paper \url{https://arxiv.org/abs/2510.04618}
\\
- MAESTRO paper \url{https://arxiv.org/abs/2601.00481}

\end{document}
